\lyseis{Β3.5}{Ιδιότητες ισοσκελούς και ισόπλευρου τριγώνου}
\sol{1} α) Σωστό. β) Σωστό: η διάμεσος προς τη βάση είναι και ύψος και διχοτόμος. γ) Σωστό: σε ισόπλευρο κάθε διάμεσος είναι και ύψος, και περνά από το μέσο της πλευράς κάθετα, άρα είναι μεσοκάθετός της. δ) Λάθος: ένας άξονας έχει το ισοσκελές που δεν είναι ισόπλευρο, τρεις το ισόπλευρο. ε) Σωστό.
\sol{2} Το $\mathrm{Μ}$ είναι μέσο της $\mathrm{ΒΓ}$ (στο ισοσκελές το ύψος προς τη βάση είναι και διάμεσος): $\mathrm{ΒΜ}=8:2=4$ cm. Το ύψος είναι κάθετο: $\mathrm{ΑΜ}\hat{\mathrm{Β}}=90^\circ$. Είναι και διχοτόμος: $\mathrm{ΒΑΜ}=50^\circ:2=25^\circ$.
\sol{3} α) $\hat{\Gamma}=\hat{\mathrm{Β}}=70^\circ$, γιατί οι γωνίες της βάσης ισοσκελούς είναι ίσες. β) $\hat{\mathrm{Β}}=\hat{\Gamma}=48^\circ$, για τον ίδιο λόγο.
\sol{4} Όλες οι πλευρές είναι ίσες: $\mathrm{ΒΓ}=6$ cm. Η $\mathrm{ΑΔ}$ είναι διάμεσος: $\mathrm{ΒΔ}=6:2=3$ cm. Είναι και ύψος: $\mathrm{ΑΔ}\hat{\mathrm{Β}}=90^\circ$. Είναι και διχοτόμος: $\mathrm{ΒΑΔ}=60^\circ:2=30^\circ$.
\sol{5} Ελεύθερη κατασκευή. Οι δύο γωνίες της βάσης βγαίνουν ίσες (περίπου $53^\circ$ η καθεμία). Ο άξονας συμμετρίας περνά από το $\mathrm{Α}$ και από το μέσο της $\mathrm{ΒΓ}$.
\sol{6} Ελεύθερη κατασκευή. Οι τρεις άξονες περνούν από το ίδιο σημείο και κάθε γωνία βγαίνει $60^\circ$.
\sol{7} α) Έναν. β) Τρεις. γ) Κανέναν.
\sol{8} Οι γωνίες της βάσης είναι ίσες, άρα και η άλλη είναι $35^\circ$.
\sol{9} Η διάμεσος προς τη βάση χωρίζει τη βάση στη μέση: $12:2=6$ cm από κάθε άκρο. Είναι και ύψος, άρα σχηματίζει με τη βάση γωνία $90^\circ$.
\sol{10} Η διχοτόμος: $\mathrm{ΒΑΔ}=64^\circ:2=32^\circ$. Η διχοτόμος προς τη βάση είναι και ύψος, άρα $\mathrm{ΑΔ}\hat{\mathrm{Β}}=90^\circ$.
\sol{11} Περίμετρος $3\cdot40=120$ cm. Έχει τρεις άξονες συμμετρίας.
\sol{12} Δεν έχει απαραίτητα δίκιο. Αν οι δύο μισές συμπίπτουν, το τρίγωνο έχει άξονα συμμετρίας, άρα είναι τουλάχιστον ισοσκελές. Είναι ισόπλευρο μόνο αν διπλώνεται με τρεις διαφορετικούς τρόπους.
\sol{13} Όταν διπλώνουμε κατά τον άξονα, το $\mathrm{Β}$ πέφτει πάνω στο $\Gamma$, άρα και η γωνία $\hat{\mathrm{Β}}$ πέφτει πάνω στην $\hat{\Gamma}$: οι δύο γωνίες εφαρμόζουν, άρα είναι ίσες.
\sol{14} Επειδή $\mathrm{ΑΒ}=\mathrm{ΑΓ}$, το $\mathrm{Α}$ ισαπέχει από τα $\mathrm{Β}$ και $\Gamma$, άρα βρίσκεται πάνω στη μεσοκάθετο της $\mathrm{ΒΓ}$. Και το $\mathrm{Μ}$, ως μέσο της $\mathrm{ΒΓ}$, ανήκει σε αυτή. Η ευθεία που περνά από τα $\mathrm{Α}$ και $\mathrm{Μ}$ είναι η μεσοκάθετος.
\sol{15} Ελεύθερη κατασκευή. Η διάμεσος προς την $\mathrm{ΑΓ}$ δεν είναι ούτε ύψος ούτε διχοτόμος, στη γενική περίπτωση (αυτό ισχύει μόνο για τη διάμεσο προς τη βάση).
\sol{16} Ελεύθερη κατασκευή. Οι δύο άλλες γωνίες βγαίνουν ίσες, περίπου $45^\circ$ η καθεμία.
\sol{17} Ελεύθερη κατασκευή. Οι δύο γωνίες της βάσης βγαίνουν ίσες.
\sol{18} Ελεύθερη κατασκευή. Οι τρεις άξονες περνούν από το ίδιο σημείο, μέσα στο τρίγωνο.
